% Part: normal-modal-logic
% Chapter: filtrations
% Section: introduction

\documentclass[../../../include/open-logic-section]{subfiles}

\begin{document}

\olfileid[fr]{nml}{fil}{int}

\olsection{Introduction}

Une question importante à propos d'une logique est toujours celle de
sa décidabilité : existe-t-il une procédure effective qui réponde à
la question « cette !!{formula} est-elle valide ? » La logique
propositionnelle est décidable : pour déterminer si une !!{formula}
est une tautologie, on peut construire sa table de vérité, laquelle
est finie pour toute !!{formula} donnée. Mais il n'est pas évident de
tester si une formule modale est vraie dans tous les modèles, car
il en existe une infinité. Pour chaque taille finie, nous pouvons
énumérer les modèles pertinents pour une formule donnée : seules
comptent les affectations des ensembles de mondes aux
!!{propositional variable}s qui figurent effectivement dans cette
!!{formula}. Une fois fixée la relation d'accessibilité, les valeurs
possibles de $V(p)$ sont les sous-ensembles de~$W$ ; si $\card{W} =
n$, il y en a $2^n$. Si notre !!{formula}~$!A$ contient $m$
!!{propositional variable}s, il existe donc $2^{nm}$ valuations
différentes sur~$n$ mondes. Pour chacune de ces valuations, on peut tester
si~$!A$ est vraie en tous les mondes en calculant la valeur de vérité
de~$!A$ en chacun d'eux. Il faut aussi examiner les relations d'accessibilité possibles,
mais leur nombre sur~$n$ mondes est fini : c'est le nombre de
sous-ensembles de $W \times W$, soit $2^{n^2}$.

Si nous étudions non pas \Log{K}, mais une logique définie par une
classe de modèles, par exemple les modèles réflexifs et transitifs,
il faut également pouvoir tester si la relation d'accessibilité est
du type requis. C'est possible lorsque les cadres considérés sont
définissables par une famille finie de !!{formula}s modales : on peut, par exemple,
tester si \Ax{T} et \Ax{4} sont valides dans un cadre fini. L'idée
serait donc de parcourir tous les cadres finis, de vérifier pour
chacun s'il appartient à la classe voulue, puis d'énumérer tous les
modèles possibles sur ce cadre et de tester si $!A$ est vraie dans
chacun. Si ce n'est pas le cas, on s'arrête : $!A$ n'est pas valide
dans la classe de modèles considérée.

Cette idée présente une difficulté : nous ne savons pas quand nous
pourrons cesser de chercher, ni même si cela arrivera. Si la
!!{formula} possède un contre-modèle fini, la procédure le trouvera.
Si elle n'en possède pas, nous n'obtiendrons pas de réponse. Cette !!{formula} peut
être valide, sans aucun contre-modèle, ou n'avoir que des
contre-modèles infinis, que nous n'examinerons jamais. Pour résoudre
cette difficulté, il suffirait de montrer que toute !!{formula} ayant
un contre-modèle en a un fini. Une logique qui possède cette propriété
a la \emph{propriété du modèle fini}.

Comment établir la propriété du modèle fini ? Une méthode consiste à
transformer un (contre-)modèle infini de~$!A$ en un modèle fini. Si
cette transformation est possible, tout modèle où $!A$ est fausse
donne un modèle fini où $!A$ est également fausse. Ce dernier
figurera dans notre énumération, qui finira donc par déceler la
non-validité de toute formule non valide. En revanche, la procédure ne
s'arrêtera pas si la formule est valide. Si nous obtenons en plus une
borne effective sur la taille du modèle fini produit, ne dépendant
que de la !!{formula}~$!A$, nous saurons jusqu'à quelle taille
examiner les modèles finis pour chercher un contre-modèle. Si aucun
n'apparaît avant cette borne, il n'en existe pas. La procédure pourra
alors répondre, pour toute !!{formula}~$!A$, à la question : « $!A$ est-elle valide ? »

Pour transformer une structure infinie en une structure finie, on
peut souvent « identifier » les !!{element}s qui se comportent de la
même manière sous les aspects pertinents. Si une infinité de mondes
de~$\mModel{M}$ ont le même comportement pertinent, nous pouvons
peut-être répartir \emph{tous} les mondes en un nombre fini de classes
d'équivalence, dont certaines peuvent être infinies. Autrement dit,
il faut définir une partition adéquate de l'ensemble des mondes :
elle comporte un nombre fini de classes, chacune pouvant être finie
ou infinie. On définit alors un nouveau modèle~$\mModel{M^*}$ dont
les mondes sont ces classes. Leur nombre fini donne un modèle fini.
La classe du monde~$w$ sera notée $[w]$. Nous voudrons avoir
$\mSat{M}{!A}[w]$ si et seulement si
$\mSat{M^*}{!A}[{[w]}]$, afin que le nouveau modèle soit un
contre-modèle de~$!A$ dès que l'ancien l'est. Il faudra pour cela
définir convenablement la partition et la relation d'accessibilité
de~$\mModel{M^*}$.

Pour comprendre le mécanisme, supposons d'abord que l'accessibilité
soit universelle : tout monde accède à tout monde. Alors
$\mSat{M}{\Box !B}[w]$ si et seulement si, pour tout $v \in W$,
$\mSat{M}{!B}[v]$ ; il en va de même pour $\mModel{M^*}$, en
remplaçant les mondes par leurs classes $[w]$ et $[v]$. Comme première
idée, disons que deux mondes $u$ et $v$ sont équivalents, c'est-à-dire
appartiennent à la même classe, s'ils s'accordent sur toutes les
!!{propositional variable}s de~$\mModel{M}$ pertinentes pour la formule étudiée :
$\mSat{M}{p}[u]$ si et seulement si $\mSat{M}{p}[v]$ pour chacune de
ces variables. Posons, pour toute variable propositionnelle,
$V^*(p) = \Setabs{[w]}{\mSat{M}{p}[w]}$. Nous voulons montrer que
$\mSat{M}{!A}[w]$ si et seulement si
$\mSat{M^*}{!A}[{[w]}]$. Une récurrence convient. Dans le cas de
base, $!A \equiv p$. Supposons d'abord $\mSat{M}{p}[w]$. Par
définition, $[w] \in V^*(p)$ ; donc $\mSat{M^*}{p}[{[w]}]$.
Réciproquement, si $\mSat{M^*}{p}[{[w]}]$, alors $[w] \in V^*(p)$ :
il existe donc un $v$ équivalent à~$w$ tel que $\mSat{M}{p}[v]$.
Dire que $w$ est équivalent à $v$ signifie que $w$ et $v$ rendent vraies les
mêmes !!{propositional variable}s pertinentes ; d'où
$\mSat{M}{p}[w]$. Pour l'étape de récurrence, considérons par exemple
$!A \ident \lnot !B$. On a $\mSat{M}{\lnot !B}[w]$ si et seulement si
$\mSat/{M}{!B}[w]$, si et seulement si
$\mSat/{M^*}{!B}[{[w]}]$ par hypothèse de récurrence, si et seulement
si $\mSat{M^*}{\lnot !B}[{[w]}]$. Les autres connecteurs non modaux
se traitent de même. Le cas de $\Box$ fonctionne également : si
$\mSat{M^*}{\Box !B}[{[w]}]$, alors, pour tout $[u]$,
$\mSat{M^*}{!B}[{[u]}]$. L'hypothèse de récurrence donne, pour tout
$u$, $\mSat{M}{!B}[u]$ ; donc $\mSat{M}{\Box !B}[w]$. La
réciproque se montre de la même manière.

Dans le cas général, où il faut également définir la relation
d'accessibilité du modèle~$\mModel{M^*}$, la situation est plus
délicate. Nous appellerons \emph{filtration} un modèle~$\mModel{M^*}$
dont la relation d'accessibilité~$R^*$ satisfait les conditions
nécessaires à la récurrence précédente. Toute filtration~$\mModel{M^*}$
rend $!A$ vraie en~$[w]$ si et seulement si $\mModel{M}$ rend
$!A$ vraie en~$w$. Il faut
toutefois prouver que de telles filtrations \emph{existent}, autrement
dit que l'on peut définir~$R^*$ de façon à satisfaire ces conditions.
Pour cela, les mondes $u$ et $v$ doivent être équivalents non seulement
lorsqu'ils s'accordent sur les !!{propositional variable}s, mais
lorsqu'ils s'accordent sur toutes les sous-!!{formula}s de~$!A$.
Comme $!A$ n'a qu'un nombre fini de sous-!!{formula}s, la filtration
reste finie. Il existe plusieurs façons de définir une filtration ;
pour que sa relation d'accessibilité ait les propriétés requises,
par exemple la réflexivité ou la transitivité, il faudra choisir
judicieusement la définition de~$R^*$.
\begin{explain}
Pour une formule à $m$ variables et une taille $n$ fixée,
on examine $2^{n^2}$ relations et, pour chacune,
$2^{nm}$ valuations sur les variables pertinentes. Une borne
effective sur la taille des contre-modèles transforme cette
énumération en un test fini. Dans l'exemple où l'accessibilité
est universelle, tout monde accède à tous les autres, ainsi
qu'à lui-même. La vérité d'une nécessité ne dépend donc pas
du monde où on l'évalue.
\end{explain}
\end{document}
